\cotefile{157-3}
\scanpage{207}
\mspage{156}

du foncteur $\varphi_Q\colon\mathcal X\longrightarrow Y$, qu'on a \unsure{oublié} d'expliciter dans 3.2. Ça a vraiment l'air de marcher (comme sur des roulettes).

Appliquons le lemme 2 \deleted{\unsure{au}} au diagramme carré $Q_{/x_0}$ et à $\varphi_{(Q_{/x_0})}$ (page 154), on voit donc que les fibres à examiner \inserted{en les $\deleted{\unsure{x|u}}u\in\operatorname{Ob}Y_{/y}$ $(u\colon y'\longrightarrow y)$} peuvent s'écrire comme des intégrales

\diagram{\begin{tikzcd}[ampersand replacement=\&,row sep=small,column sep=small]
\& (X_{1/x_1})_u\\
\mathcal X(x_0,u)\overset{\text{déf}}{=}\displaystyle\int (X_{0/x_0})_u \arrow[ur]\arrow[dr] \&\\
\& (X_{2/x_2})_u .
\end{tikzcd}}

\deleted{\unsure{U}}Utilisons maintenant le fait que $f_1,f_2$ sont $W$-lisses (exprimant que $\varphi_Q$ est $W$-lisse dans $\mathcal X-X_0$), donc \deleted{que} \unsure{:} $(X_{1/x_1})_u$ et $(X_{2/x_2})_u$ sont $W$-asphériques. Dans le carré $W$-cocartésien

\diagram{\begin{tikzcd}[ampersand replacement=\&,row sep=small,column sep=small]
\& (X_{1/x_1})_u \arrow[dr,"{\unsure{f_1|i_1|j_1}(x_0,u)}",""{name=branchabove,near start}] \&\\
(X_{0/x_0})_u \arrow[ur,"{\alpha_1(x_0,u)}"]\arrow[dr,"{\alpha_2(x_0,u)}"'] \& \& (\mathcal X_{/x_0})_u \simeq \mathcal X(x_0,u) ,\\
\& (X_{2/x_2})_u \arrow[ur,"{i_2(x_0,u)}"',""{name=branchbelow,near start}]\arrow[from=branchabove,to=branchbelow,dashed,no head,shorten <=4pt,shorten >=4pt] \&
\end{tikzcd}}

% slanted
\margin{$\circ$ Attention \unsure{ce|on} n'est pas un carré commutatif. Dire plutôt que $(\mathcal X_{/x_0})_u\simeq\int$
\diagram{\begin{tikzcd}[ampersand replacement=\&,row sep=0pt,column sep=4pt]
\& {}\\[-3pt]
{} \arrow[ur]\arrow[dr] \&\\[-3pt]
\& {}
\end{tikzcd}}
est lisse sur $\underline{\Phi}$, à fibres $W$-asphériques (si $f_0,f_1,f_2$ sont lisses) donc $W$-as, \unsure{il faudrait que puisque} $\underline{\Phi}$ \unsure{l'est}.}

% slanted
\margin{\deleted{Au lieu de $W$\unsure{$(7')$} (bis) il suffit de \unsure{de plus|plus} supposer que \unsure{Cat muni du système} \unsure{est une 2-cat.} ? \unsure{$(\underline W,\underline{\Phi W})$} \unsure{cofinal triang.} ! \unsure{?} pour la \unsure{partie} \unsure{$f_1,f_2,f_0$ lisses} $\Rightarrow$ \unsure{$\varphi$ lisse}. C'est \unsure{moins fort} que cette \unsure{assertion}, car alors les \unsure{diff.} iso, sont \unsure{des} \illegible{1}.}}

``$(\mathcal X_{/x_0})_u$ est $W$-\unsure{asphérique|asph.}'' \unsure{et} $(X_{1/x_1})_u$ et $(X_{2/x_2})_u$ le \unsure{sont|sont-elles} \unsure{équivalent|équivalant} à $i_1(x_0,u)\in W$ (ou $i_2(x_0,u)\in W$) donc aussi \unsure{(moyennant $W(7\ \mathrm{bis})$)} :
% archive sheet number: 207
